\section{极值格截面与设计矩}
\label{sec:sections}

43维计算与五维截面正交的极小向量数；45维将三维截面的三个纤维投影到
正交补。完整投影域与设计矩能够确定所需统计量，无须枚举整个母壳。

\subsection{母壳与矩方程}
48维极值偶幺模格的最小平方范数为六，theta 级数为
\[
E_4^6-1440E_4^3\Delta+125280\Delta^2
=1+52416000q^3+O(q^4).
\]
这里 $E_4=1+240q+2160q^2+\cdots$，
$\Delta=q-24q^2+252q^3+\cdots$；常数项及前两个零系数确定这个权24模形式
\cite{conway1999sphere}。极小壳 $X$ 是球面11-design
\cite[定理4.1]{nebe-designs}。

对 Gram 矩阵为 $H$ 的截面向量 $t_i$，令
$y_i=x\cdot t_i$，$f(y)=|\{x\in X:y(x)=y\}|$。
整数性、$y^{\mathsf T}H^{-1}y\le6$ 及与截面短向量的内积界限制签名域。
当 $2m\le10$ 时，
\begin{equation}\label{eq:section-moment}
\sum_{x\in X}\prod_{\ell=1}^{2m}(x\cdot t_{i_\ell})
=\frac{|X|6^m}{48\cdot50\cdots(48+2m-2)}
\sum_{\mathcal P}\prod_{(a,b)\in\mathcal P}H_{i_ai_b},
\end{equation}
其中 $\mathcal P$ 遍历带标号位置的配对；次数零给出总点数。
这里包含总次数为偶数、个别指数为奇数的单项式。
若矩系统为 $Av=b$，有理恒等式 $A^{\mathsf T}w=c$ 即证明
$c^{\mathsf T}v=w^{\mathsf T}b$。平均必须保持完整域、球面矩和目标，
但不要求延拓为母格自同构。

\subsection{43维：固定子截面与八个坐标截面}
在 $P_{48n}$ 中~\cite{latticecatalogue}，给定整数行矩阵 $T$ 嵌入
$\sqrt3E_8$：母格 Gram 矩阵 $G$ 满足
$(TGT^{\mathsf T})_{ij}=3\alpha_i\cdot\alpha_j$，其中
\[
\alpha_1=\tfrac12(1,-1,-1,-1,-1,-1,-1,1),\quad
\alpha_2=e_1+e_2,\quad
\alpha_j=e_{j-1}-e_{j-2}\quad(3\le j\le8).
\]
前五行指定 $\sqrt3D_5$ 子截面 $U$。嵌入母截面的路线属于
Dorofeev--Sun--Wang 的截面框架~\cite{dorofeev}。

\begin{proposition}[$E_8$ 扩张计数]\label{prop:section43}
设 $L$ 是48维极值偶幺模格，含等距嵌入
$\iota:\sqrt3E_8\hookrightarrow L$，$X_L$ 为其平方范数六的壳。
上述前五个嵌入简单根的张成空间 $U_L$，以及这个子截面的每个 $W(E_8)$ 共轭，
都满足
\[
|X_L\cap U_L^\perp|=2553792.
\]
不要求母格具有相应 Weyl 作用，也不要求嵌入本原。
给定嵌入由此实现 $K(43)\ge2553792$。
\end{proposition}

由整数性与自对偶性，母截面投影写成 $\lambda/\sqrt3$，
其中 $\lambda\in E_8$。除去240个单点例外 $\lambda=3r$，
最小范数条件给出 $\|\lambda\|^2\le18$，且对每个根 $r$ 有
$|\lambda\cdot r|\le3$。精确枚举在平方范数 $0,2,4,6,8$ 上分别得到
$1,240,2160,6720,17280$ 个签名，加上例外共26641个。
径向矩确定 Weyl 平均重数
\[
(a_0,a_2,a_4,a_6,a_8)=(1092000,116235,9180,495,63/4).
\]
子截面的正交空间在这些轨道中分别含 $1,12,6,24,0$ 个签名，
另有十二个例外，所以平均点数为
$1092000+12(116235)+6(9180)+24(495)+12=2553792$。

为证明固定子截面的结论，取保持目标的群
$\langle W(D_5),W(A_3),-I\rangle$；其中 $A_3$ 由与子截面正交的根组成。
完整域分为43个轨道。重建矩行与例外值固定行后，证书满足
$A^{\mathsf T}w=c$、$w^{\mathsf T}b=2553792$，从而证明命题，
无需母格对称性。完整域不可省略：正交条件应放入目标 $c$，
不能预先删除其他列。不同极小向量内积至多为三，故除以 $\sqrt6$
得到接吻构型。比较构造由等距的 $D_5$ 截面给出2545056点；
改进来自嵌入位置，而非 Gram 型。

将 $\iota$ 线性延拓，定义
$U_k=\iota(\sqrt3\operatorname{span}\{e_1,\ldots,e_k\})$，
$1\le k\le8$。令 $m=8-k$，则八个固定截面的精确点数为
\begin{align*}
|X_L\cap U_k^\perp|={}&1092000+18360m+464944\binom m2
+11880\binom m3\\
&+146880\binom m4+2520\binom m5+31680\binom m6.
\end{align*}
这只陈述指定坐标截面及其 Weyl 共轭的点数，不涉及任意截面的最优性。
$U_1$ 是坐标轴而非 $D_1$ 根系；其余是坐标 $D_k$ 的张成空间。

具体地，在倍坐标 $u=2\lambda$ 中，按前 $k$ 个及后 $8-k$ 个坐标内的
置换与偶数次变号对完整非例外域分轨。使用
$m_a(u_1^2,\ldots,u_k^2)m_b(u_{k+1}^2,\ldots,u_8^2)$，
其中 $m_a$ 为单项对称多项式、$|a|+|b|\le5$；
相应矩由 Gram 矩阵 $12I_8$ 给出。八个轨道系统的满列秩依次为
$25,34,39,41,39,34,25,13$，由非零整数子式认证。
因此实际轨道总重数唯一，等于 Weyl 平均值 $|O|a_r$。
再计算目标签名数即得公式；只用径向矩则只能得到截面间的平均。
Weyl 变换
$\tfrac12P_{5,8}\operatorname{diag}(J_4-2I_4,J_4-2I_4)$
将指定 $D_5$ 的张成空间送到前五条坐标轴，其中 $P_{5,8}$ 交换第五与第八坐标。

\begin{center}
\begin{tabular}{r|rrrr}
$k$&1&2&3&4\\\hline
$|X_L\cap U_k^\perp|$&16815624&10663920&6688960&4149504\\[2pt]
$k$&5&6&7&8\\\hline
$|X_L\cap U_k^\perp|$&2553792&1593664&1110360&1092000
\end{tabular}
\end{center}

\subsection{45维：精确纤维恒等式}
附录~\ref{app:ambient}证明所指定二次剩余偶邻格 $\Lambda$ 的最小范数为六。
以下计数更一般地适用于任意含 Gram 矩阵
$H=8I_3-2J_3$、$H^{-1}=(I_3+J_3)/8$ 截面的48维极值偶幺模格，
不要求截面本原。

截面向量 $t_1,t_2,t_3,-t_1-t_2-t_3$ 及其负向量给出八个单点签名
\[
E=\{\pm He_1,\pm He_2,\pm He_3,\pm H(1,1,1)^{\mathsf T}\}.
\]
其余签名均属于
\begin{equation}\label{eq:domain45}
D_0=\{y\in\mathbb Z^3:|y_i|\le3,\ |y_1+y_2+y_3|\le3\}.
\end{equation}
这231个签名的投影平方范数至多为 $9/2$。
扣除 $E$ 并合并反极点后，矩系统有161行、116个变量，
$A_{a,[y]}=(2-\mathbf1_{y=0})y^a$，$b_a=M_a-\sum_{y\in E}y^a$。

\begin{proposition}[精确纤维恒等式]\label{prop:moment45}
在上述条件下，
\[
f(e_1)+f(e_2)+f(e_3)=7377408+9f(-2,-2,3).
\]
\end{proposition}
\begin{proof}
令 $c$ 在 $[e_i]$ 上取系数1，在 $[(-2,-2,3)]$ 上取 $-9$，其余为零。
附录~\ref{app:data}的 \texttt{dual.json} 给出有理乘子，满足
$c-A^{\mathsf T}w\ge0$、$w^{\mathsf T}b=7377408$。
但 $t_i+t_j$ 的平方范数为八，最小范数又给出 $|y_i+y_j|\le4$。
写 $y_4=-\sum_{i=1}^3y_i$，这些条件恰好排除
$(3,3,-3,-3)$ 与 $(3,2,-3,-2)$ 的30个排列。
七个正松弛项全部属于这些被排除的轨道，故在剩余101个反极轨道上
$A_D^{\mathsf T}w=c_D$，原下界证书因而加强为等式。
\end{proof}

三个纤维投影后的平方范数为 $23/4$，同一纤维内积至多为 $11/4$，
跨纤维内积至多为 $23/8$。这同时排除投影重合，归一化即为接吻构型。
因此目标纤维中 $m$ 个不同见证给出 $K(45)\ge7377408+9m$；
完整纤维的大小则给出构型的精确点数。

\subsection{锚图与认证点数}
对平方范数为八的锚点 $p$，令 $A_p=\{x\in X:x\cdot p=4\}$。
最小范数给出 $|x\cdot p|\le4$。滤波多项式
$P(t)=t^2(t^2-1)(t^2-4)(t^2-9)$ 满足 $P(4)=20160$，
全壳求和为90961920，反极对称性因而给出 $|A_p|=2256$。
将 $x$ 与 $p-x$ 配对。不同点对的端点内积属于 $\{1,2,3\}$，
更换端点将内积变为四减去原值。因此内积一或三定义一个1128顶点的图。

对 $z\perp p$，次数至多十的混合矩经滤波给出
$\sum_{x\in A_p}(z\cdot x)^2=192\|z\|^2$。
取 $z=r-p/2$，自身点对贡献32，每个相邻点对贡献二，
于是 $768=32+2\deg(\bar r)$：每个这样的锚图都是368正则图。

对由 $r,s$ 代表的不相邻顶点，截面 $(-s,s-p,p-r)$ 的 Gram 矩阵为 $H$。
$N(\bar r)\setminus N(\bar s)$ 中每个顶点有唯一端点 $a$ 满足
$r\cdot a=3$，此时 $w=p-a$ 的签名为 $(-2,-2,3)$。
反之，这个签名强制 $p\cdot w=4$、$r\cdot w=1$、$s\cdot w=2$，
从而恢复唯一端点。完整纤维大小因此为 $368-c$，
其中 $c=|N(\bar r)\cap N(\bar s)|$，精确投影点数为 $7380720-9c$。
每个 Gram 为 $H$ 的有序截面都可这样表示：
$p=-t_1-t_2$、$s=-t_1$、$r=-t_1-t_2-t_3$。

每对选择一个端点 $r_i$，令 $u_i=r_i-p/2$。
带符号矩阵 $S$ 对角为零，非对角元为
$S_{ij}=u_i\cdot u_j\in\{0,\pm1\}$。上述二阶矩给出
\[
\sum_i u_i\otimes u_i=96I_{p^\perp},\qquad
S^2=88S+368I,\qquad
\operatorname{Spec}(S)=\{92^{(47)},(-4)^{(1081)}\}.
\]
这是因为 Gram 矩阵 $4I+S$ 的秩为47，平方等于自身的96倍。
带符号共同邻点数满足 $c_{ij}^{\pm}=(c_{ij}\pm88S_{ij})/2$，
所以共同邻点数均为偶数，投影点数被18整除。
无符号邻接矩阵满足 $A+16I\succeq0$，因为它是张量 $u_i\otimes u_i$
的 Gram 矩阵；这并不意味着无符号谱只有两个值。
四个带符号的 $u_i$ 不可能两两内积均为 $-1$，
否则其和是平方范数为四的格向量。

所选锚点在缩放因子 $1/\sqrt{12}$ 的分子坐标中具有形状 $(3^6,1^{42})$。
2256个认证端点组成完整母集合。全部428076对不相邻顶点的共同邻点数
介于70与148之间，唯一一对达到70，故该纤维恰有298点，
改进了此前 Pless 格的292见证构造。

\begin{corollary}[45维构造]\label{cor:section45}
指定构型恰有 $7377408+9(298)=7380090$ 点，因此 $K(45)\ge7380090$。
这是固定锚点及 Gram 矩阵 $H$ 下的最大投影点数，
不是对其他锚点、母格或截面 Gram 矩阵的全局最优性结论。
\end{corollary}
